\section{模型的评价推广与生产建议}

\subsection{模型的优点}

\begin{enumerate}[label=\arabic*., leftmargin=2.6em, itemsep=4pt]
  \item 本文没有直接套用现成的干燥模型，而是先用扩散的特征时间与表面传质阻力判断主控机制——水分扩散约需 22 h、干燥由内部扩散控制，从而说明烘干时长对表面传质系数并不敏感，也解释了预热阶段为何只有表层失水，体现了建模过程的数学严谨性与思考的深度。
  \item 本文把四个问题统一在同一套模型框架下，只在物性与环境条件上分支，问题二与问题三能互相印证；对不断内缩的边界创新性地引入归一化坐标，把移动的边界转化为一个已知的对流项，并给出可独立判错的检验，逻辑严谨。
  \item 本文绘制了许多可视化图像让论文更加得直观，易于读者的理解。
\end{enumerate}

\subsection{模型的缺点}

\begin{enumerate}[label=\arabic*., leftmargin=2.6em, itemsep=4pt]
  \item 模型按一维径向轴对称处理，只保留径向坐标，
        未计及药材两端面的轴向输运与热风沿周向的不均匀性。
  \item 数值求解的计算量较大。主方法取 $800$ 个区间，烘干过程长达数十小时，
        单次全时程计算需要推进数十万个时间步，耗时较长。

\end{enumerate}

\subsection{模型的推广}

本文的建模与求解框架不限于中药材烘干。对于内部扩散控制、表面为对流边界的圆柱或球形物料，
只要替换物性经验式与环境条件，即可用于其他干燥、吸附与浸提过程的模拟；
问题四处理收缩边界的坐标变换方法，也可用于计算域随时间变化的溶解、膨胀与烧蚀等扩散问题。

\subsection{对实际生产的建议}

热风烘干依赖烘房温湿环境的调控，本文得到的干燥规律可直接落到以下四条生产建议上。

\begin{enumerate}[label=\arabic*., leftmargin=2.6em, itemsep=4pt]
  \item \textbf{调控重点是烘房温度，不是风速。} 
  \item \textbf{排产必须计入尺寸收缩，且前期升温宜缓。} 
  \item \textbf{停机应以含水率判定并留余量。} 
  \item \textbf{按模型排产需预留约一成的系统偏差。}
\end{enumerate}
